\clearpage\fancyhead[L]{\small COMET.}\hypertarget{native-350-lower}{}\label{native-350-lower}
\begin{center}\textbf{COMET.}\par\vspace{4pt}\hrule\vspace{4pt}\fontsize{10}{12}\selectfont Göttingische gelehrte Anzeigen. Stück 15. Seite 147..149. 1812. Januar 25.\par\vspace{4pt}\hrule\vspace{4pt}\end{center}

\par\indent Die seit länger als einem Monate ungewöhnlich ungünstige Witterung hat seit den ersten, in diesen Blättern (St. 201 von vor. J. [S. 342 d. B.]) angezeigten, Beobachtungen des \emph{neuen} Cometen nur noch drei Beobachtungen auf hiesiger Sternwarte verstattet, die wir zugleich mit den erstern, in der Reduction etwas verbessert, mittheilen:

\begin{center}
\fontsize{10}{12}\selectfont\setlength{\tabcolsep}{3pt}\renewcommand{\arraystretch}{1.15}\begin{tabular}{l|c|c|c}
 & \shortstack{Mittlere\\Zeit} & \shortstack{Gerade\\Aufsteig.} & Abweichung \\
\hline
1811 Decemb. 9 & $10^{\text{u}}6^{\text{m}}52^{\text{s}}$ & $63^\circ49'41''4$ & $10^\circ21'55''5$ S. \\
11 & $10\;\;34\;\;1$ & $63\;\;33\;\;18.0$ & $8\;\;39\;\;46.4$ \\
12 & $8\;\;5\;\;52$ & $63\;\;26\;\;25.3$ & $7\;\;54\;\;25.9$ \\
1812 Januar 3 & $7\;\;1\;\;30$ & $62\;\;37\;\;6.4$ & $9\;\;47\;\;14.7$ N. \\
4 & $6\;\;53\;\;59$ & $62\;\;40\;\;38.6$ & $10\;\;29\;\;57.8$ \\
\end{tabular}
\end{center}

\noindent Schon am 13. December hatte Hr. Prof. \textsc{Gauss} aus seinen ersten Beobachtungen und denen des Hrn. v. \textsc{Zach}, welche wir gleichfalls a.\,a.\,O. mitgetheilt haben, vorläufige parabolische Elemente berechnet, die im Westphälischen Moniteur Nr.\,301 und im Decemberheft der Monatl. Correspondenz abgedruckt sind. Natürlich konnten diese auf so dürftige Data gegründeten Elemente nur genäherte sein, welche indessen am 3. Jan. nur erst 8 Min. von der Beobachtung abwichen. Hr. Prof. \textsc{Gauss} übertrug nun, nachdem er seine Beobachtung vom 4. Jan. gemacht hatte, die Verbesserung der parabolischen Elemente Hrn. \textsc{Nicolai}, welcher diese Arbeit mit grösster Sorgfalt und dem besten Erfolge ausgeführt hat. Fol\SourceLeaf{351}gendes sind die verbesserten Elemente, welche dieser geschickte Rechner herausgebracht hat, und die gewiss nur noch unbedeutender Verbesserungen fähig sein werden.

\Needspace{15\baselineskip}
\begin{center}\fontsize{10}{12}\selectfont\setlength{\tabcolsep}{4pt}\renewcommand{\arraystretch}{1.15}\begin{tabular}{l}
Zeit der Sonnennähe 1811 Nov. 11\\
\quad $9^{\mathrm u}\;24^{\mathrm m}\;38^{\mathrm s}$ Merid. von Göttingen.\\
\end{tabular}\end{center}
\begin{center}\fontsize{10}{12}\selectfont\setlength{\tabcolsep}{3pt}\renewcommand{\arraystretch}{1.15}\begin{tabular}{llr}
Abstand in der Sonnennähe & .\;.\;.\;.\;.\;.\;.\;.\;.\;.\;. & $1.589198$\\
Länge der Sonnennähe & .\;.\;.\;.\;.\;.\;.\;.\;.\;.\;.\;.\;. & $47^\circ\;39'\;36''1$\\
Länge des aufsteigenden Knoten & .\;.\;.\;.\;.\;.\;.\;.\;. & $92\;53\;44.2$\\
Neigung der Bahn & .\;.\;.\;.\;.\;.\;.\;.\;.\;.\;.\;.\;.\;. & $31\;32\;38.7$\\
\end{tabular}\end{center}
\noindent Bewegung rechtläufig.

\par\indent Die Constanten, vermittelst welcher nach Hrn. Prof. \textsc{Gauss} Methode die Coordinaten des Cometen in Beziehung auf den Aequator berechnet werden, sind folgende:

\Needspace{13\baselineskip}
\[
\begin{aligned}
x\cos\tfrac{1}{2}\upsilon^2&=\alpha\sin(\upsilon+138^\circ\;9'\;39''6)\\
y\cos\tfrac{1}{2}\upsilon^2&=\ObservedBeta\sin(\upsilon+59\;53\;53.5)\\
z\cos\tfrac{1}{2}\upsilon^2&=\gamma\sin(\upsilon+355\;26\;5.4)\\
\log\alpha&=0.1319477\\
\log\ObservedBeta&=0.1784770\\
\log\gamma&=9.9865984
\end{aligned}
\]

\Needspace{35\baselineskip}\par\indent Ausser den oben angeführten Beobachtungen sind Hrn. Prof. \textsc{Gauss} noch einige auf der Seeberger Sternwarte von Hrn. v. \textsc{Lindenau}, zu Bremen von Hrn. Dr. \textsc{Olbers}, und zu Paris von Hrn. \textsc{Burckhardt} angestellte mitgetheilt worden; letztere sind Meridianbeobachtungen. Mit diesen sämmtlichen Beobachtungen verglich Hr. \textsc{Nicolai} seine Elemente, und erhielt folgende schöne Uebereinstimmung, die kaum noch Etwas zu wünschen übrig lässt.

\begin{center}
\fontsize{10}{12}\selectfont\setlength{\tabcolsep}{5pt}\renewcommand{\arraystretch}{1.15}\begin{tabular}{l|c|c|l}
 & \multicolumn{2}{c|}{Unterschied} & \\
 & \shortstack{in gerader\\Aufsteigung} & \shortstack{in der\\Abweichung} & Beobachter \\
\hline
1811 Nov. 18 & $-12''4$ & $-138''4$ & v. Zach \\
19 & $+34.7$ & $+19.5$ & v. Zach \\
20 & $-11.9$ & $-5.5$ & v. Zach \\
21 & $-45.0$ & $+4.9$ & v. Zach \\
Dec. 8 & $+9.7$ & .\;.\;.\;. & v. Lindenau \\
9 & $+1.0$ & $-51.4$ & Olbers \\
9 & $+4.9$ & $+22.6$ & v. Lindenau \\
9 & $-5.0$ & $+3.4$ & Gauss \\
11 & $+2.9$ & $-32.7$ & Gauss \\
12 & $+8.5$ & $-27.6$ & Gauss \\
14 & $+23.5$ & $-34.4$ & Olbers \\
14 & $+18.9$ & $-6.7$ & Burckhardt \\
16 & $+28.7$ & $-18.1$ & Olbers \\
22 & $+20.8$ & $+2.4$ & Burckhardt \\
25 & $+25.1$ & $-3.1$ & Burckhardt \\
1812 Jan. 3 & $+2.7$ & $-21.1$ & Olbers \\
3 & $-5.5$ & $+12.9$ & Gauss \\
4 & $+2.8$ & $-12.7$ & Gauss \\
6 & $-35.6$ & $-1.5$ & Olbers \\
\end{tabular}
\end{center}

\par\indent Der Comet hat seit den ersten hiesigen Beobachtungen, wo seine Lichtstärke am grössten und etwa $= 0.7$ war, merklich an Licht abgenommen, doch liess er sich am 3. und 4. Januar noch eben so gut, wie vorher, beobachten. Er würde noch sehr lange sichtbar bleiben, wenn nicht zu befürchten \SourceLeaf{352}wäre, dass sein nun immer mehr abnehmendes Licht ihn wohl schon im Februar unsern Augen entziehen wird. Inzwischen werden die Beobachtungen sehr durch folgende Ephemeride erleichtert werden die gleichfalls von Hrn. \textsc{Nicolai} berechnet ist, und welcher zufolge der Comet, nachdem er bereits mitten durch die Hyaden gegangen ist, aus dem Stier in den Fuhrmann treten und dort wahrscheinlich unsichtbar werden wird.
